% Style follows ../class2/resume2PT26.tex.
% Build in this directory (run twice for page references):
% pdflatex -interaction=nonstopmode -halt-on-error resume11PT26.tex
% Topics follow the syllabus in ../class1/class1PT26.tex.
% Content adapted from ../lectureNote.tex and ../lectureNotes.tex.
\documentclass[a4paper,10pt,pdflatex,ja=standard,jafont=ipaex-type1,
  magstyle=real,scale=1]{bxjsarticle}
\usepackage{lmodern}
\usepackage{amsmath,amssymb,mathtools}
\usepackage{geometry}
\geometry{reset,left=22mm,right=22mm,top=20mm,bottom=17mm,
  headheight=13pt,headsep=6mm,footskip=10mm}
\usepackage{xcolor}
\usepackage{enumitem}
\usepackage{array,booktabs,tabularx}
\usepackage{tikz}
\usepackage{fancyhdr}
\usepackage{hyperref}
\definecolor{ink}{HTML}{193A5C}
\definecolor{accent}{HTML}{237D82}
\definecolor{soft}{HTML}{EAF3F3}
\definecolor{muted}{HTML}{52616B}
\newcommand{\coursename}{確率論と統計学 I}
\hypersetup{unicode=true,colorlinks=true,urlcolor=accent,linkcolor=ink,
  pdfauthor={中島秀太},pdftitle={第11回：条件付き期待値}}
\pagestyle{fancy}
\fancyhf{}
\fancyhead[L]{\small\color{ink}\coursename{}・授業要約と演習}
\fancyhead[R]{\small\color{ink}第11回}
\fancyfoot[C]{\small\color{ink}\thepage/\pageref*{lastpage}}
\renewcommand{\headrulewidth}{0.3pt}
\setlength{\parindent}{0pt}
\setlength{\parskip}{2pt}
\setlength{\abovedisplayskip}{4pt}
\setlength{\belowdisplayskip}{4pt}
\setlength{\abovedisplayshortskip}{3pt}
\setlength{\belowdisplayshortskip}{4pt}
\setlist[itemize]{leftmargin=1.8em,itemsep=1pt,topsep=2pt,parsep=0pt}
\setlist[enumerate]{leftmargin=2em,itemsep=3pt,topsep=3pt,parsep=0pt}
\newcommand{\topic}[1]{\par\vspace{5pt}
  {\large\sffamily\mdseries\color{accent}#1}\par\vspace{2pt}}
\newcommand{\handouttitle}[1]{%
  \begin{center}
    {\Large\color{ink}#1}\par\vspace{3mm}
    {\color{accent}\rule{0.88\linewidth}{0.7pt}}
  \end{center}\vspace{-2mm}}
\newcommand{\keybox}[1]{\par\vspace{3pt}\begingroup
  \setlength{\fboxsep}{5pt}
  \colorbox{soft}{\parbox{\dimexpr\linewidth-2\fboxsep\relax}{#1}}%
  \endgroup\par\vspace{2pt}}
\newcommand{\examplelabel}[1]{{\sffamily\mdseries\color{ink}#1}}
\newcommand{\columnheading}[1]{\par\vspace{5pt}
  {\fontsize{11pt}{15pt}\selectfont\sffamily\mdseries\color{accent}#1}\par\vspace{2pt}}
\newcommand{\columnstart}[1]{\clearpage
  \setlength{\parskip}{2.5pt}
  \fancyhead[L]{\small\color{ink}\coursename{}・発展的内容}
  {\small\sffamily\mdseries\color{accent}発展的内容}\par\vspace{4pt}
  {\Large\color{ink}#1}\par\vspace{2pt}
  {\color{accent}\rule{\linewidth}{0.7pt}}\par\vspace{4pt}}
\newcommand{\F}{\mathcal F}
\newcommand{\R}{\mathbb R}
\newcommand{\N}{\mathbb N}
\newcommand{\E}{\mathbb E}
\newcommand{\Var}{\operatorname{Var}}
\newcommand{\Cov}{\operatorname{Cov}}
\newcommand{\ind}{\mathrm{1}}
\newcommand{\iid}{\text{独立同分布}}
\begin{document}
\fontsize{10pt}{13.5pt}\selectfont
\begingroup
\fancyhead[L]{\small\color{ink}\coursename{}・前回の演習解答}
\handouttitle{第10回\quad 演習問題の解答}

\topic{演習1：指数分布の標本平均}
独立同分布な $X_j\sim\mathrm{Exp}(\lambda)$（$\lambda>0$）と
$\bar X_n=n^{-1}\sum_{j=1}^nX_j$ について
$$
  \E[\bar X_n]=\frac1\lambda,\qquad
  \Var(\bar X_n)=\frac1{n\lambda^2},\qquad
  P\left(\left|\bar X_n-\frac1\lambda\right|>\varepsilon\right)
     \le\frac1{n\lambda^2\varepsilon^2}\quad(\varepsilon>0).
$$
最後の評価は Chebyshev の不等式による。

\topic{演習2：中心化した独立な和の四次モーメント}
$Y_j$ は独立同分布で $\E[Y_j]=0$、$\mu_4=\E[Y_j^4]<\infty$、
$v=\E[Y_j^2]$ とする。四乗を展開した積で、ある添字が一度だけ現れるなら、
独立性によってその $\E[Y_j]=0$ を因子としてもつので期待値は消える。
残るのは同一の添字4個、または異なる添字が2個ずつの場合だけである。
$$
  \E\left[\left(\sum_{j=1}^nY_j\right)^4\right]
     =n\mu_4+6\binom n2v^2=n\mu_4+3n(n-1)v^2.
$$
係数 $6$ は、固定した二つの添字を2個ずつ並べる $\binom42$ 通りから生じる。

\topic{演習3：指数関数を与える極限}
$a_n\to a\in\R$ なら $(a_n)$ は有界であり、十分大きい $n$ で
$|a_n/n|\le1/2$、従って $1-a_n/n>0$ となる。
$\log(1-u)=-u+O(u^2)$（$u\to0$）を使うと
$$
  n\log(1-a_n/n)=-a_n+O(a_n^2/n)\longrightarrow-a.
$$
指数関数の連続性より、$(1-a_n/n)^n\to e^{-a}$ を得る。

\topic{演習4：サイコロの大数の法則と中心極限定理}
公平なサイコロを独立に投げた出目を $D_1,D_2,\ldots$ とし、
$\bar D_n=n^{-1}\sum_{j=1}^nD_j$ とおく。
$$
  \E[D_j]=\frac72,\qquad
  \E[D_j^2]=\frac{91}{6},\qquad
  \Var(D_j)=\frac{91}{6}-\frac{49}{4}=\frac{35}{12}.
$$
有界なので強大数の法則から $\bar D_n\to7/2$ a.s.、従って確率収束もする。
中心極限定理は
$$
  \sqrt{\frac{12n}{35}}\left(\bar D_n-\frac72\right)
      \Rightarrow N(0,1)
$$
を与える。前者は平均の安定化、後者は $n^{-1/2}$ の尺度で見た揺らぎの分布を述べる。
\clearpage
\endgroup

\handouttitle{第11回\quad 条件付き期待値}
\keybox{条件付き期待値は、手元の情報だけから決まる確率変数である。
可測性と、情報で区別できるすべての事象上での積分一致が、その定義となる。}

\topic{有限空間での条件付き分布}
$(\Omega,\F,P)$ を有限確率空間、$X,Y$ を実数値確率変数とする。
$P(Y=y)>0$ に対して
$$P(X=x\mid Y=y)=\frac{P(X=x,Y=y)}{P(Y=y)},\qquad
 f(y):=\E[X\mid Y=y]=\sum_x xP(X=x\mid Y=y).$$
$P(Y=y)=0$ の値では $f(y)=0$ など任意に定めてよい。
この関数に実際の観測 $Y$ を代入して
$$\E[X\mid Y]:=f(Y)$$
と定める。$\E[X\mid Y=y]$ は数、$\E[X\mid Y]$ は $\Omega$ 上の確率変数である。

\topic{例：二回のコイントス}
公平なコインを独立に2回投げ、$X$ を1回目が表なら $1$、裏なら $0$、
$Y$ を表の総数とする。
\begin{center}
\renewcommand{\arraystretch}{1.25}
\begin{tabular}{@{}ccccc@{}}
\toprule
$\omega$ & $HH$ & $HT$ & $TH$ & $TT$ \\
\midrule
$X$ & $1$ & $1$ & $0$ & $0$ \\
$Y$ & $2$ & $1$ & $1$ & $0$ \\
$\E[X\mid Y]$ & $1$ & $1/2$ & $1/2$ & $0$ \\
\bottomrule
\end{tabular}
\end{center}
$Y=1$ と知ったとき、$HT,TH$ は等確率なので $\E[X\mid Y=1]=1/2$。
この例では $\E[X\mid Y]=Y/2$ となる。

\topic{情報を表す $\sigma$-加法族}
$\mathcal G=\sigma(Y)=\{Y^{-1}(B):B\in\mathcal B(\R)\}$ は、
$Y$ の観測だけで判定できる事象の族である。
有限空間では、同じ $Y$ の値をもつ標本を一つのまとまりに分けて考えられる。
一般の部分 $\sigma$-加法族 $\mathcal G\subseteq\F$ も、利用可能な情報を表す。

\topic{一般の定義：二つの条件}
$\E|X|<\infty$ とする。確率変数 $Z$ が
$$Z\text{ は }\mathcal G\text{-可測},\qquad \E|Z|<\infty,\qquad
 \E[Z\ind_A]=\E[X\ind_A]\quad(\forall A\in\mathcal G)$$
を満たすとき、$Z$ を $\E[X\mid\mathcal G]$ と書く。
可測性は $\mathcal G$ の情報だけで値を決められる条件、積分一致は
その情報で区別できるどの事象上でも平均が保たれる条件である。
$\E[X\mid Y]$ は $\E[X\mid\sigma(Y)]$ の略記として使う。

\newpage
\topic{存在と一意性}
$A\in\mathcal G$ に対して $\nu(A)=\E[X\ind_A]$ と定める。
$\nu$ は有限符号付き測度で $\nu\ll P|_{\mathcal G}$ なので、Radon--Nikodym の定理が
$$Z=\frac{d\nu}{d(P|_{\mathcal G})}$$
という可積分な $\mathcal G$-可測関数を与える。これが存在を保証する。
同じ定義を満たす $Z_1,Z_2$ には、$A=\{Z_1>Z_2\}\in\mathcal G$ を使うと
$\E[(Z_1-Z_2)\ind_A]=0$ なので $P(Z_1>Z_2)=0$。
逆の不等式も同様で、$Z_1=Z_2$ a.s. となる。

\topic{線形性と、すでに分かっている量}
可積分な $X,Y$ と $a,b\in\R$ に対して
$$\E[aX+bY\mid\mathcal G]
 =a\E[X\mid\mathcal G]+b\E[Y\mid\mathcal G]\quad\text{a.s.}$$
$X$ が $\mathcal G$-可測なら $\E[X\mid\mathcal G]=X$ a.s.。
また有界な $\mathcal G$-可測確率変数 $H$ について
$$\E[HX\mid\mathcal G]=H\E[X\mid\mathcal G]\quad\text{a.s.}$$
となる。既知の量は条件付き期待値の外へ出せる。
有界性を外す場合には、積の可積分性を別に確認する必要がある。

\topic{全期待値と塔の性質}
定義式に $A=\Omega$ を入れるだけで
$$\E[\E[X\mid\mathcal G]]=\E[X]$$
を得る。さらに $\mathcal H\subseteq\mathcal G\subseteq\F$ なら
$$\E[\E[X\mid\mathcal G]\mid\mathcal H]
 =\E[X\mid\mathcal H]\quad\text{a.s.}$$
である。左辺は $\mathcal H$-可測で、任意の $A\in\mathcal H$ に対し積分をとると
$$\E[\E[\E[X\mid\mathcal G]\mid\mathcal H]\ind_A]
 =\E[\E[X\mid\mathcal G]\ind_A]=\E[X\ind_A].$$
定義と一意性から結論が従う。

\topic{情報がない場合と独立な情報の場合}
自明な $\sigma$-加法族 $\{\varnothing,\Omega\}$ については
$\E[X\mid\{\varnothing,\Omega\}]=\E[X]$ a.s.。
$X$ が $\mathcal G$ と独立、すなわち $\sigma(X)$ と $\mathcal G$ の任意の事象が独立なら
$$\E[X\mid\mathcal G]=\E[X]\quad\text{a.s.}$$
となる。実際、定数 $\E[X]$ は可測であり、
$\E[X\ind_A]=\E[X]P(A)$ が積分一致を与える。

\keybox{候補を計算したら、情報についての可測性と積分一致を確かめる。
等式は原則として a.s. の意味であり、零確率の条件上の値は一意に定まらない。}

\newpage
\topic{演習問題}
\begin{enumerate}[label=\arabic*.]
\item 公平なコインを独立に3回投げ、$X$ を1回目の表の指示関数、$T$ を表の総数とする。
$\E[X\mid T=t]$（$t=0,1,2,3$）を求め、$\E[X\mid T]$ を $T$ の式で表せ。
\item 公平なサイコロの出目を $D$ とし、$\mathcal G$ を偶奇だけを知る情報とする。
$\E[D\mid\mathcal G]$ の値とその分布を求め、全期待値の公式を確認せよ。
\item $X$ が可積分で $\mathcal G$-可測なら、$X$ 自身が条件付き期待値の定義を満たすことを示せ。
さらに $\E[\E[X\mid\mathcal G]]=\E[X]$ を定義から示せ。
\item $\mathcal H\subseteq\mathcal G$ とする。塔の性質を定義から証明せよ。
情報の包含関係を逆にすると、同じ式をそのまま使えるか考えよ。
\end{enumerate}

\topic{レポート問題}
\begin{enumerate}[label=\arabic*.]
\item 公平なサイコロ2個を独立に投げ、出目を $U,V$、$S=U+V$ とする。
$\E[U\mid S]$ を求め、その確率分布を明示せよ。
対称性から候補を考えた後、各 $S=s$ の条件付き分布で確認すること。
\item $\E[X^2]<\infty$ とし、$Z=\E[X\mid\mathcal G]$ とおく。
任意の二乗可積分な $\mathcal G$-可測変数 $H$ に対して
$$\E[(X-H)^2]=\E[(X-Z)^2]+\E[(Z-H)^2]$$
を示せ。$Z$ の二乗可積分性には条件付き Jensen の不等式を使ってよい。
\end{enumerate}

\topic{考えるための手がかり}
条件付き期待値は、観測で区別できない標本の間を平均する。
有限空間では、情報による分割を図や表にすると計算しやすい。
一般の定義の証明では、候補が可測であることも忘れず確認しよう。
\vfill
{\small\color{muted}解答メモ：数 $\E[X\mid Y=y]$ と確率変数 $\E[X\mid Y]$ を区別する。}

\columnstart{条件付き期待値は、情報に基づく最良予測}
見えている情報だけを使って、まだ分からない量 $X$ を予測したい。
予測誤差を二乗で測るとき、条件付き期待値には明確な最適性がある。
この性質は「情報を $\sigma$-加法族で表す」という考え方を幾何学へつなぐ。

\columnheading{予測の候補は $\mathcal G$-可測な変数}
$\E[X^2]<\infty$、$Z=\E[X\mid\mathcal G]$ とする。
条件付き Jensen の不等式より
$$Z^2\le\E[X^2\mid\mathcal G]\quad\text{a.s.},\qquad
 \E[Z^2]\le\E[X^2]<\infty.$$
情報 $\mathcal G$ だけを使う予測値 $H$ も、$\mathcal G$-可測な二乗可積分変数とする。
$X-Z$ は、情報から説明できずに残る誤差である。

\columnheading{残った誤差は既知の量と直交する}
定義の積分一致を有界な $\mathcal G$-可測関数へ拡張すれば
$\E[(X-Z)K]=0$ となる。
二乗可積分な $K$ へは切断と Cauchy--Schwarz の不等式で拡張できる。
特に $K=Z-H$ と置くと
$$\E[(X-Z)(Z-H)]=0.$$
したがって二乗を展開した交差項が消え、
$$\E[(X-H)^2]=\E[(X-Z)^2]+\E[(Z-H)^2]
 \ge\E[(X-Z)^2].$$
最小値は $H=Z$ で達成され、等号なら $H=Z$ a.s.。
これは $L^2$ 空間での直交射影に対応している。

\columnheading{情報が増えると、最小誤差は増えない}
$\mathcal H\subseteq\mathcal G$ なら、$\mathcal H$-可測な予測はすべて
$\mathcal G$-可測でもある。選べる予測が増えるので
$$\E[(X-\E[X\mid\mathcal G])^2]
 \le\E[(X-\E[X\mid\mathcal H])^2].$$
コイン2回の例で、何も知らなければ1回目の表の指示関数の予測は $1/2$。
表の総数を知れば $Y/2$ となり、1回目の結果を直接知れば誤差は $0$ になる。
情報の包含関係が、予測精度の関係として表れる。

\columnheading{時間とともに更新する予測}
情報が $\F_s\subseteq\F_t$（$s\le t$）と増えるとき、
$M_t=\E[X\mid\F_t]$ とおけば塔の性質から
$$\E[M_t\mid\F_s]=M_s.$$
この性質をもつ過程を martingale という。
将来の予測値を、現在の情報で平均し直すと現在の予測に戻る。
条件付き期待値は一回の計算だけでなく、情報が増える過程の整合性も記述している。

\label{lastpage}
\end{document}
